\subsection{14.26: restricted wreath products in a nilpotent quasivariety}
\label{cand:14.26}
\problemstatement{14.26}
Let \(\mathcal C\) be all torsion-free nilpotent groups, with no fixed
class bound. If \(G\in Q(\mathcal C)\), then \(G\wr\Z\in Q(\mathcal C)\)
for the restricted wreath product. This is the direct wreath
product of the question: Budkin's definition uses finite-support
functions \cite[Section~1]{budkin-wreath}.

A finite set \(B\ni1\) locally embeds in a group if an injective
identity-preserving map preserves every product staying in \(B\).
Since \(\mathcal C\) is closed under finite direct products,
\[
 G\in Q(\mathcal C)
 \quad\Longleftrightarrow\quad
 \text{every finite }B\subseteq G,\ 1\in B,
 \text{ locally embeds in some }C\in\mathcal C.
\]
Indeed, write down the finite multiplication table of \(B\).
For each unequal pair, some realization in \(\mathcal C\) must separate
it, or the corresponding quasiidentity would be violated by \(G\).
The product of these finitely many realizations separates all pairs.
Conversely, a quasiidentity violation involves only finitely many
intermediate word values, which a local embedding would preserve.

We construct the needed local models. Fix \(c,R\ge1\) and a finite
nonempty set \(I\subseteq\Z\) of size \(M\), with diameter at most \(R\).
Choose \(m\ge2R(c-1)+M-1\), a rational vector space with basis
\(a_0,\ldots,a_m\), and a nilpotent derivation determined by
\(D(a_j)=a_{j+1}\), \(D(a_m)=0\). Put
\[
 f(z)=\sum_{j=0}^m z^ja_j/j!.
\]
In its symmetric algebra let \(J\) be the ideal generated by all
coefficients in \(z\) of \(f(z)f(z+d)\), \(1\le d\le R\), and form
\[
 \mathcal A=\operatorname{Sym}\langle a_0,\ldots,a_m\rangle/
       (J+\operatorname{Sym}^{>c}).
\]
This is finite-dimensional, commutative, and graded by ordinary degree.
The ideal is \(D\)-stable because
\(D(f(z)f(z+d))=\frac{d}{dz}(f(z)f(z+d))\).
On the degree-truncated algebra, \(D\) raises the total index of each
monomial, bounded by \(cm\), so \(D^{cm+1}=0\), also after quotienting.
Thus \(\alpha=\exp D\) is an automorphism, with
\[
 \alpha^sf(i)=f(i+s),\qquad
 f(i)f(j)=0\quad(0<|i-j|\le R).                    \tag{a}
\]

For each \(1\le k\le c\), the vectors \(f(i)^k\), \(i\in I\), are
linearly independent in \(\mathcal A\).
To prove it, a linear functional on \(\operatorname{Sym}^k\)
corresponds to a symmetric polynomial
\(Q(z_1,\ldots,z_k)=\ell(f(z_1)\cdots f(z_k))\) of degree at most
\(m\) in each variable. Conversely the coefficients of any such
polynomial prescribe \(\ell\), with factorial denominators.
If \(Q\) vanishes when \(z_2=z_1+d\), \(1\le d\le R\), the functional
kills the degree-\(k\) part of \(J\). Coefficient extraction and
Vandermonde spanning by the other \(f(z_j)\)'s justify this for every
quadratic generator multiplied by an arbitrary degree-\(k-2\) term.
Define
\[
 \Delta_k=\prod_{a<b}\prod_{d=1}^R((z_a-z_b)^2-d^2),\qquad
 C_k=\prod_{a<b}\prod_{d=1}^R(-d^2).
\]
If \(L_i\) is the degree-\((M-1)\) Lagrange polynomial on \(I\), then
\[
 Q_i=\frac{\Delta_k}{C_k}
 L_i\left(\frac{z_1+\cdots+z_k}{k}\right)
\]
has the required degree bound in each variable, kills those relations,
and has \(Q_i(j,\ldots,j)=\delta_{ij}\).
Empty products handle \(k=1\); the cutoff adds no relation in these
degrees. This proves independence.

Let \(\mathcal A_+\) be the positive-degree ideal, \(U\) the
upper unitriangular group of matrix size \((c+1)\times(c+1)\) with strict
entries in \(\mathcal A_+\), and
\[
 H=U\rtimes\langle t\rangle,\qquad tut^{-1}=\alpha(u).
\]
For \(g\in\mathrm{UT}_{c+1}(\Q)\), define \(\rho_i(g)\) by the
off-diagonal entries
\[
 \rho_i(g)_{ab}=g_{ab}f(i)^{b-a}\quad(a<b).
\]
Superdiagonal distances add, so this is a homomorphism, and the
independence above makes it injective for \(i\in I\).
Equation (a) gives \(\alpha\rho_i=\rho_{i+1}\).
For distinct positions in \(I\), the corresponding strictly upper
matrices multiply to zero in either order. Thus
\[
 \prod_{i\in I}\rho_i(g_i)=1+\sum_{i\in I}(\rho_i(g_i)-1)
\]
and the independence of each \(f(i)^k\) makes this direct-product map
injective.

The group \(H\) is torsion-free nilpotent.
On each successive superdiagonal quotient of \(U\), conjugation by
\(U\) is trivial and \(t\) acts as \(\alpha\).
Writing \(\delta=\alpha-1\), we have \(\delta^{cm+1}=0\).
Insert the images of successive powers of \(\delta\) in every
superdiagonal quotient. Their inverse images give a central series,
of length at most \(c(cm+1)+1\) including \(H/U\).
Torsion projects trivially to \(H/U=\Z\); in \(U\), a positive
power multiplies the first nonzero superdiagonal by a positive integer,
so cannot kill it over \(\Q\).

Now take a finite subset \(S\ni1\) of \(N\wr\Z\), with
\(N\in\mathcal C\).
The finitely many base coordinates generate a finitely generated
torsion-free nilpotent group, which embeds in
\(\mathrm{UT}_{c+1}(\Q)\) for some \(c\).
This classical embedding input follows from Jennings'
augmentation-separation theorem \cite[Theorem~5.2]{jennings}:
a finite-dimensional quotient \(\Q[N]/\Delta^s\) is faithful,
and the \(\Delta\)-filtration makes the action unitriangular.
Choose \(I\) to contain all supports in \(S\) and all shifted supports
needed in products \(uv=w\) within \(S\); take \(I=\{0\}\) if empty.
With \(R\) bounding its diameter, the map
\[
 (f,n)\longmapsto\left(\prod_i\rho_i(f(i))\right)t^n
\]
preserves those products: conjugation by \(t^n\) shifts the indices,
different positions commute, and equal positions multiply correctly.
It is injective by the quotient \(H/U\) and the direct-product
independence above. Hence every such \(S\) locally embeds in \(\mathcal C\).

Finally, if \(G\in Q(\mathcal C)\), collect from a finite
\(S\subseteq G\wr\Z\) all base coordinates and all coordinate products
needed by its finite multiplication table.
Locally embed this finite set in a group \(N\in\mathcal C\).
Coordinate replacement, leaving the integer exponents unchanged,
locally embeds \(S\) into \(N\wr\Z\), where the preceding construction
applies. The local criterion proves \(G\wr\Z\in Q(\mathcal C)\).
The model and its nilpotency class depend on \(S\); no global
embedding of the wreath product in a nilpotent group is asserted.
See \artifact{research/14.26-proof.md}.
